\section{Conclusion}

In this study, a real-time, explainable crowd anomaly detection system was developed, designed to identify early-stage panic and crowd crushes in surveillance video. Addressing the limitations of existing optical-flow pipelines, three distinct modifications were introduced: (1)~the expansion of the motion descriptor from three to nine features to capture angular coordination and velocity dispersion, (2)~the replacement of the unidirectional temporal architecture with a Bidirectional LSTM coupled with an additive attention mechanism, and (3)~the integration of feature-level and temporal-level post-hoc interpretability modules.

The results of our study validate the proposed system. On the UMN benchmark dataset, the BiLSTM+Attention classifier got 99.55\% accuracy and 99.27\% F1-Score, showing 100\% precision (no false positives in the test data). The warning lead time evaluation showed that anomaly events were detected 0.39~s before their peak on average, with the maximum lead time of 7.5~s. The proposed system also runs at 54.5~FPS on a standard laptop GPU, which is sufficient for real-time edge processing. Finally, the zero-shot experiments on the UCSD Ped2 and CUHK Avenue datasets confirm that our feature descriptor can generalize to different layouts, resolutions, and camera views without retraining.

However, several limitations remain. The system exhibits sensitivity to severe image degradation, with classification metrics dropping significantly under simulated sensor noise and motion blur. Additionally, the dense optical flow formulation assumes a static camera perspective, and the pipeline has not been evaluated on dynamic pan-tilt-zoom (PTZ) camera systems or live, non-laboratory CCTV streams.

Future research will focus on addressing these boundary conditions. First, self-supervised pretraining and adversarial training strategies will be explored to improve feature robustness under low-light and high-noise conditions. Second, global motion compensation techniques will be integrated into the preprocessing stage to enable compatibility with PTZ camera systems. Third, the deployment of the pipeline on resource-constrained edge hardware will be evaluated under real-world network transmission constraints to transition the system from a laboratory setting to active operational environments.
